\chapter{Real VLA란 무엇인가}

\section*{학습 목표}
이 장의 목표는 VLA를 ``큰 VLM에 action head를 붙인 모델''로 이해하는 관점을 정리하고, 실제 로봇 시스템에서 필요한 인터페이스를 명확히 분리하는 것이다. 이 장을 마치면 독자는 다음 질문에 답할 수 있어야 한다.

\begin{enumerate}
  \item VLM, VLA policy, Real VLA system의 차이를 설명할 수 있다.
  \item 정책 수준 action $\act_t$와 제어 입력 $\ctrlcmd_t$를 구분할 수 있다.
  \item 언어 명령이 실제 로봇 명령으로 바뀌는 경로를 layer 단위로 분해할 수 있다.
  \item VLA 논문을 backbone, action representation, data, control interface, evaluation, safety 기준으로 읽을 수 있다.
\end{enumerate}

\section{VLA는 더 큰 VLM이 아니다}

Vision-language model(VLM)은 이미지나 비디오와 텍스트 사이의 의미적 대응을 학습한다. 예를 들어 VLM은 장면 속 물체를 설명하거나, 이미지에 대한 질문에 답하거나, 자연어 지시를 문장 수준에서 해석할 수 있다. 그러나 로봇에서 필요한 것은 텍스트 답변이 아니라 물리적 실행이다. 로봇은 관측한 장면과 주어진 언어 지시를 동역학, 접촉, 지연, 센서 오차, actuator limit, 안전 제약을 가진 물리 행동으로 바꾸어야 한다.

따라서 VLA를 단순히 ``VLM 뒤에 action head를 붙인 모델''이라고 정의하면 중요한 부분이 사라진다. Action head가 어떤 형식의 행동을 내는지, 그 행동이 어떤 controller로 들어가는지, 안전 layer가 무엇을 검사하는지, 실패했을 때 누가 멈추고 누가 복구하는지, 실험 결과가 어떤 evidence로 기록되는지가 모두 Real VLA의 일부이다.

\begin{definitionbox}{핵심 정의: Real VLA}
이 책에서 Real VLA는 특정 모델 이름이 아니다. Real VLA는 language-conditioned perception과 learned action을 실제 로봇의 물리 실행, monitoring, evaluation, safety assurance로 연결하는 책임 구조이다.
\end{definitionbox}

이 관점은 고전 로봇공학을 대체하지 않는다. 오히려 고전 로봇공학의 개념을 더 선명하게 만든다. 안정성, 제약 만족, 접촉 제어, 모션 feasibility, state estimation, failure recovery는 VLM의 파라미터 수가 커져도 사라지지 않는다. VLA가 실제 로봇에 가까워질수록 이 문제들은 더 중요해진다.

\subsection{정책 action과 제어 입력}

이 장에서 가장 먼저 고정할 notation은 정책 수준 action과 제어 입력의 차이이다. VLA policy가 출력하는 action을 $\act_t$라 쓰고, 실제 servo loop 또는 robot driver가 받는 제어 입력을 $\ctrlcmd_t$라 쓰자. 일반적으로
\begin{equation}
  \act_t = \policy(\obs_t, \conf_t, \task, \bel_t),
  \label{eq:vla-policy}
\end{equation}
이다. 여기서 $\obs_t$는 sensor observation, $\conf_t$는 robot configuration, $\task$는 언어 지시 또는 subgoal, $\bel_t$는 state belief 또는 memory를 뜻한다. 그러나 $\act_t$는 actuator command가 아니다. 로봇에 실제로 들어가는 명령은 보통 다음 변환을 거친다.
\begin{equation}
  \ctrlcmd_t = \controller(\act_t, \hat{\state}_t, \mathcal{C}_t),
  \label{eq:controller}
\end{equation}
여기서 $\hat{\state}_t$는 추정 상태이고, $\mathcal{C}_t$는 kinematic constraint, collision constraint, force limit, workspace bound, safety condition 등을 포함한다.

\begin{warningbox}{주의: VLA는 controller replacement가 아니다}
VLA policy는 grasp, waypoint, delta pose, action token, action chunk를 제안할 수 있다. 하지만 그 자체가 collision-free trajectory, stable contact, actuator-safe torque, distribution shift recovery를 보장하지 않는다. Real VLA engineering의 큰 부분은 $\act_t \rightarrow \ctrlcmd_t$ 경계에 있다.
\end{warningbox}

\section{언어 목표에서 모터 명령까지}

다음 명령을 보자.
\begin{quote}
머그컵을 drying rack 위에 올려라. 단, 옆의 유리컵은 건드리지 마라.
\end{quote}

이 문장은 간단해 보이지만 실제 로봇 시스템에서는 여러 종류의 정보로 분해되어야 한다.

\subsection{Task semantics}
첫 번째 변환은 language에서 task semantics로 가는 변환이다. ``머그컵''은 조작 대상이고, ``drying rack''은 목표 위치 또는 목표 relation이며, ``유리컵을 건드리지 말라''는 hard constraint이다. ``올려라''라는 동사는 grasp, lift, transport, place, release, verify로 이어지는 작업 구조를 암시한다. 즉 자연어 지시는 action 자체가 아니라 action을 조직하는 task-level specification이다.

\subsection{Embodied reasoning}
두 번째 변환은 task semantics에서 embodied reasoning으로 가는 변환이다. 시스템은 머그컵과 유리컵을 찾고, pose 또는 affordance를 추정하고, drying rack의 가능한 placement region을 계산하며, 유리컵 주변 forbidden region을 설정해야 한다. 이 단계는 LLM, VLM, embodied reasoning model, scene graph, memory, planner, tool call을 사용할 수 있다. 그러나 출력은 여전히 motor command가 아니다.

\subsection{Policy-level action}
세 번째 변환은 embodied reasoning 결과를 policy-level action으로 바꾸는 것이다. VLA policy는 observation, language 또는 subgoal, proprioception, short memory를 받아 다음 action object를 출력한다. 이 action object는 end-effector delta, waypoint, discrete action token, short action chunk, continuous action distribution 등일 수 있다. 이 출력이 $\act_t$이다.

\subsection{Control command}
네 번째 변환은 $\act_t$를 controller target 또는 low-level command로 바꾸는 것이다. 예를 들어 VLA가 end-effector delta pose를 냈다면 motion/control layer는 현재 robot state, coordinate frame, kinematic limit, obstacle geometry, controller rate를 고려하여 이를 해석해야 한다. Operational-space controller, impedance controller, MPC, trajectory optimizer, whole-body controller 등이 여기서 $\ctrlcmd_t$를 만든다.

\begin{figure}[ht]
\centering
\begin{tikzpicture}[
  node distance=8mm,
  block/.style={draw=rvlaBlue, rounded corners=1mm, fill=rvlaGray, align=center, minimum width=32mm, minimum height=9mm, font=\sffamily\small},
  arrow/.style={-Latex, thick, draw=rvlaBlue}
]
\node[block] (inst) {언어 지시\\task context};
\node[block, below=of inst] (er) {Embodied reasoning\\scene, constraint, subgoal};
\node[block, below=of er] (vla) {VLA policy\\$\obs_t,\conf_t,\task,\bel_t \mapsto \act_t$};
\node[block, below=of vla] (ctrl) {Motion/control\\$\act_t,\hat{\state}_t,\mathcal{C}_t \mapsto \ctrlcmd_t$};
\node[block, below=of ctrl] (robot) {Robot/environment\\execution and feedback};
\node[block, right=10mm of vla, minimum width=29mm] (safe) {Safety/assurance\\allow, modify, stop};
\node[block, left=10mm of vla, minimum width=29mm] (data) {Data/evaluation\\logs, metrics, failures};
\draw[arrow] (inst) -- (er);
\draw[arrow] (er) -- (vla);
\draw[arrow] (vla) -- (ctrl);
\draw[arrow] (ctrl) -- (robot);
\draw[arrow] (robot.west) to[out=180,in=-90] (data.south);
\draw[arrow] (data.north) to[out=90,in=180] (er.west);
\draw[arrow] (safe.west) -- (vla.east);
\draw[arrow] (safe.south west) -- (ctrl.east);
\end{tikzpicture}
\caption{Real VLA system stack. 이 그림은 특정 구현 pipeline이 아니라 책임 구조를 나타낸다.}
\label{fig:rvla-stack}
\end{figure}

\section{다섯 개의 책임 layer}

Real VLA를 시스템으로 읽기 위해 이 책은 다섯 개의 layer를 사용한다. 중요한 것은 layer가 반드시 별도 module이어야 한다는 뜻이 아니라, 책임이 사라지지 않아야 한다는 뜻이다.

\begin{table}[ht]
\centering
\caption{Real VLA stack의 입력, 출력, 책임}
\label{tab:stack}
\small
\begin{tabularx}{\textwidth}{p{24mm}X X X}
\toprule
Layer & 핵심 질문 & 대표 입력 & 대표 출력 \\
\midrule
Embodied reasoning & 무엇을 해야 하는가? & instruction, scene, memory, task context & subgoal, plan, constraint, tool/policy call \\
VLA policy & 지금 어떤 action을 제안할 것인가? & observation, language/subgoal, proprioception & policy-level action $\act_t$ \\
Motion/control & action을 어떻게 물리 motion으로 바꿀 것인가? & $\act_t$, robot state, constraints & controller command $\ctrlcmd_t$ \\
Safety/assurance & 허용, 수정, 정지, 복구 중 무엇을 할 것인가? & instruction, plan, action, runtime signals & allow, modify, stop, recover, handoff \\
Data/evaluation & 어떤 evidence가 생성되었는가? & logs, actions, states, failures, interventions & dataset record, metric, regression test \\
\bottomrule
\end{tabularx}
\end{table}

\subsection{Embodied reasoning layer}
Embodied reasoning layer는 intent를 구조화한다. 이 layer는 scene graph, memory, planner, VLM/LLM, skill library를 사용할 수 있다. 하지만 좋은 plan은 physical execution의 필요조건일 수는 있어도 충분조건은 아니다. ``유리컵을 피해서 머그컵을 옮겨라''라는 subgoal이 있다고 해서 자동으로 collision-free trajectory와 stable grasp가 생기지는 않는다.

\subsection{VLA policy layer}
VLA policy layer는 이 책에서 ``VLA''라는 단어가 가장 좁은 의미로 쓰이는 위치이다. 이 layer는 observation, language 또는 subgoal, robot state를 받아 policy-level action object를 낸다. 이 action object는 action type, coordinate frame, rate, horizon, scaling, downstream controller를 요구하는 interface object이다.

\subsection{Motion/control layer}
Motion/control layer는 단순 executor가 아니다. 같은 ``오른쪽으로 이동''이라는 action도 base frame 기준인지 end-effector frame 기준인지, normalized vector인지 metric pose인지, velocity command인지 target pose인지에 따라 전혀 다르게 실행된다. 따라서 controller는 learned policy의 하위 부품이 아니라 physical feasibility를 해석하는 별도 책임 layer이다.

\subsection{Safety/assurance layer}
Safety는 마지막 wrapper가 아니다. 위험한 instruction을 거절하는 semantic safety, forbidden region을 검사하는 geometric safety, 속도와 힘 제한을 다루는 dynamic safety, 실행 중 이상을 감지하는 runtime assurance, 사람의 intervention을 가능하게 하는 handoff mechanism이 모두 필요하다. Runtime safety decision을
\begin{equation}
  r_t = \sigma(\task, \hat{\state}_t, \act_t, \ctrlcmd_t, h_t)
  \in \{\textsf{allow},\textsf{modify},\textsf{stop},\textsf{ask}\}
  \label{eq:safety-decision}
\end{equation}
로 쓸 수 있다. 여기서 $h_t$는 최근 sensor history와 system log이다.

\subsection{Data/evaluation layer}
Data/evaluation layer는 deployment 이후 문서화 절차가 아니다. VLA의 성능 주장은 데이터와 평가 protocol 위에 선다. 어떤 robot, 어떤 camera, 어떤 gripper, 어떤 action representation, 어떤 controller, 어떤 randomization, 어떤 reset/intervention policy를 썼는지에 따라 같은 success rate의 의미가 달라진다.

\section{인터페이스 계약: observation, state, action, control}

Real VLA stack은 interface가 명확할 때만 분석 가능하다. 첫 번째 구분은 observation과 state이다. $\obs_t$는 camera image, depth, tactile signal, proprioception 등 system이 관측한 값이다. $\state_t$는 실제 환경과 로봇의 물리 상태이다. 실제 상태는 대개 직접 관측되지 않는다. 따라서 system은 state estimate $\hat{\state}_t$ 또는 belief $\bel_t$ 위에서 작동한다.

\begin{equation}
  \bel_t = \eta(\bel_{t-1}, \obs_t, \ctrlcmd_{t-1}),
  \label{eq:belief-update}
\end{equation}
와 같이 belief update를 쓸 수 있다. 여기서 $\eta$는 estimator, filter, learned state representation, memory update를 포괄한다.

두 번째 구분은 action representation이다. VLA action은 다음처럼 다양한 형태를 가질 수 있다.

\begin{table}[ht]
\centering
\caption{대표 action representation과 system-level implication}
\label{tab:action-rep}
\small
\begin{tabularx}{\textwidth}{p{28mm}X X}
\toprule
Action type & 장점 & 주의할 점 \\
\midrule
Action token & VLM/LLM sequence modeling과 잘 맞음 & token의 물리 의미와 controller mapping이 불명확할 수 있음 \\
Continuous delta & controller와 직접 연결하기 쉬움 & scaling, normalization, coordinate frame에 민감함 \\
Waypoint & planning/control integration이 쉬움 & downstream planner와 collision checking 부담이 커짐 \\
Action chunk & temporal consistency를 줄 수 있음 & interruption과 recovery granularity가 거칠어질 수 있음 \\
Diffusion/flow output & multimodal action distribution 표현 가능 & sampling latency, safety projection이 필요함 \\
\bottomrule
\end{tabularx}
\end{table}

Action representation은 formatting 문제가 아니다. 무엇을 $\act_t$로 정의하는지에 따라 policy가 학습하는 대상, controller가 해석해야 할 대상, safety layer가 검사할 수 있는 대상, evaluator가 기록해야 할 대상이 모두 달라진다.

\section{역사적 수렴: 제어, 계획, 로봇러닝, VLM/LLM}

Real VLA는 갑자기 등장한 분야가 아니다. 네 흐름이 수렴한 결과로 보는 것이 정확하다.

\begin{enumerate}
  \item \textbf{Control and estimation}: state, feedback, stability, contact, latency, constraints.
  \item \textbf{Planning and TAMP}: task structure, feasibility, long-horizon composition, recovery.
  \item \textbf{Robot learning}: visuomotor policy, imitation, RL, dataset coverage, distribution shift.
  \item \textbf{VLM/LLM/foundation models}: semantic grounding, language-conditioned reasoning, tool use.
\end{enumerate}

Control과 estimation은 VLA가 물리계에 들어갈 때 필요한 언어이다. VLA가 end-effector delta pose를 내든, waypoint를 내든, action chunk를 내든, 로봇은 여전히 kinematics, dynamics, contact, actuator limit 안에서 움직인다.

Planning과 TAMP도 현재형 도구이다. 자연어 instruction은 종종 long-horizon structure를 갖는다. LLM 또는 VLM이 planning function 일부를 흡수할 수는 있지만, planning problem 자체가 사라지는 것은 아니다. 문제는 PDDL을 반드시 쓰느냐가 아니라 task structure, constraint, feasibility, recovery가 system 어딘가에 존재하느냐이다.

Robot learning은 VLA의 직접 조상이다. VLA 이전에도 behavior cloning, DAgger, guided policy search, visual RL, large-scale grasping은 image-to-action mapping을 학습하려고 했다. VLA가 새롭게 가져온 것은 여기에 language grounding과 foundation-model representation을 결합했다는 점이다.

\section{현대 VLA system pattern}

현대 VLA system은 다음 pattern으로 읽으면 좋다. 이것은 완성된 taxonomy가 아니라 논문을 읽기 위한 engineering map이다.

\begin{table}[ht]
\centering
\caption{현대 VLA system pattern}
\label{tab:patterns}
\small
\begin{tabularx}{\textwidth}{p{28mm}X X}
\toprule
Pattern & 핵심 기여 & 과대해석 위험 \\
\midrule
VLA policy & vision/language/proprioception에서 action을 직접 예측 & policy를 full robot system으로 오해함 \\
Action-distribution model & token, chunk, diffusion, flow 등 action 생성 구조를 설계 & controller와 safety assumption이 숨겨짐 \\
ER + VLA orchestration & reasoning model이 task decomposition 또는 VLA/skill call 수행 & reasoning competence를 motor competence로 오해함 \\
Humanoid/whole-body stack & loco-manipulation, balance, whole-body control까지 확장 & demo를 broad deployment evidence로 오해함 \\
Open-source training stack & fine-tuning, action tokenizer, benchmark harness 제공 & local deployment가 해결되었다고 오해함 \\
\bottomrule
\end{tabularx}
\end{table}

이 pattern들은 모두 같은 결론으로 이어진다. VLA 연구의 핵심은 backbone 크기만이 아니다. 어떤 action representation을 쓰는지, 어떤 controller에 연결되는지, 어떤 data record가 남는지, 어떤 evaluation protocol이 claim을 지지하는지가 함께 중요하다.

\section{Data, evaluation, safety는 appendices가 아니다}

Real VLA claim은 architecture만으로 성립하지 않는다. Policy가 image와 language를 받아 action을 낼 수 있어도, 실제 capability는 data regime, evaluation protocol, safety layer에 의존한다.

\subsection{Robot data의 구조}
Robot trajectory는 단순한 observation sequence가 아니다. 다음 record를 포함해야 해석 가능하다.
\begin{equation}
\begin{aligned}
\mathcal{D}_i
&=
(\task_i,O_i,Q_i,A_i,U_i,m_i,y_i,e_i),\\
O_i &= \{\obs_t\}_{t=1}^{T_i},
\qquad
Q_i = \{\conf_t\}_{t=1}^{T_i},\\
A_i &= \{\act_t\}_{t=1}^{T_i},
\qquad
U_i = \{\ctrlcmd_t\}_{t=1}^{T_i}.
\end{aligned}
\label{eq:data-record}
\end{equation}
여기서 $m_i$는 robot, controller, action type, rate 같은 metadata이고, $y_i$는 success, failure, intervention label이며, $e_i$는 safety event, reset policy, normalization 기록이다. 이 정보가 없으면 같은 demonstration도 다르게 해석된다. 5 Hz Cartesian delta action을 compliant controller가 tracking한 trajectory와 50 Hz joint target을 position controller가 tracking한 trajectory는 같은 training signal이 아니다.

\subsection{Benchmark success와 robot capability}
Benchmark는 필요하다. 공통 task, 비교 가능한 score, regression test를 제공하기 때문이다. 그러나 benchmark success가 곧 robot capability는 아니다. Benchmark는 일부 변수를 고정하여 특정 차원을 측정한다. Deployment는 여러 차원을 동시에 노출한다.

머그컵 예제에서 Real VLA evaluation은 단순히 ``머그컵이 rack 위에 올라갔는가''만 묻지 않는다. 최소한 다음 항목을 봐야 한다.
\begin{itemize}
  \item mug가 목표 region에 놓였는가?
  \item glass와 contact가 없었는가?
  \item contact force와 velocity limit이 유지되었는가?
  \item human intervention이 있었는가?
  \item slip이나 실패를 감지하고 recover했는가?
  \item layout, lighting, distractor 변화에서 재현되는가?
  \item $\act_t$와 $\ctrlcmd_t$의 interface assumption이 기록되었는가?
\end{itemize}

Evaluation vector를 다음처럼 정의할 수 있다.
\begin{equation}
  \bm{m} =
  \left[
  p_{\mathrm{succ}},
  p_{\mathrm{safe}},
  \tau_{\mathrm{latency}},
  n_{\mathrm{intervention}},
  r_{\mathrm{recovery}},
  \Delta_{\mathrm{shift}}
  \right].
  \label{eq:evaluation-vector}
\end{equation}
여기서 $p_{\mathrm{succ}}$만 보고 capability를 주장하면 stack의 중요한 failure mode가 가려진다.

\begin{warningbox}{Benchmark success는 robot capability가 아니다}
Benchmark result는 system이 특정 protocol을 만족했다는 뜻이다. Robot capability는 더 넓은 claim이다. Protocol, randomization, failure taxonomy, intervention policy, safety event, controller assumption이 보이지 않으면 number의 의미도 제한된다.
\end{warningbox}

\subsection{Safety as stack property}
Safety는 polite language model이나 collision checker 하나로 환원되지 않는다. Real VLA safety는 semantic, geometric, dynamic, runtime, data, human-supervision 책임을 포함한다.

\begin{table}[ht]
\centering
\caption{Safety responsibility by layer}
\label{tab:safety}
\small
\begin{tabularx}{\textwidth}{p{30mm}X X}
\toprule
Layer & Safety question & Example failure \\
\midrule
Semantic & instruction이 허용 가능한가? & constraint를 preference로 해석함 \\
Geometric & 경로와 목표가 충돌 없이 가능한가? & glass 근처 forbidden region을 무시함 \\
Dynamic & 속도, 힘, torque가 제한 안에 있는가? & place 동작 중 overshoot 발생 \\
Runtime & 실패를 충분히 빨리 감지하는가? & contact 후에야 stop signal 발생 \\
Data & unsafe data가 학습/평가에 섞였는가? & hidden reset이 success로 기록됨 \\
Human & 사람이 중단하거나 인계받을 수 있는가? & intervention path가 없음 \\
\bottomrule
\end{tabularx}
\end{table}

\section{VLA 논문을 읽는 여섯 가지 질문}

이 책의 나머지 장에서는 다음 reading protocol을 반복적으로 사용한다. 어떤 VLA 논문이 robot capability를 주장할 때, 먼저 다음 여섯 질문을 던져야 한다.

\begin{enumerate}
  \item \textbf{Backbone은 무엇인가?} VLM, LLM, diffusion model, transformer policy, flow model, hierarchical planner, hybrid architecture 중 무엇인가?
  \item \textbf{Action representation은 무엇인가?} token, Cartesian delta, pose, waypoint, action chunk, joint target, torque command, latent action 중 무엇을 출력하는가?
  \item \textbf{Data는 무엇으로 구성되는가?} real robot trajectory, simulation, teleoperation, scripted demonstration, human video, synthetic data, web pretraining 중 무엇인가?
  \item \textbf{Control interface는 무엇인가?} policy output을 어떤 controller가 받는가? control frequency와 tracking assumption은 명시되었는가?
  \item \textbf{Evaluation protocol은 무엇인가?} trial count, randomization, distribution shift, failure taxonomy, intervention policy가 보이는가?
  \item \textbf{Safety/deployment evidence는 무엇인가?} instruction filtering, geometric check, runtime monitor, stop/recovery, human override, audit log가 있는가?
\end{enumerate}

이 질문들은 단순한 비기술적 요약표가 아니다. 어떤 interface가 바뀌었고 어떤 evidence가 추가되었는지 분리하기 위한 engineering rubric이다. 어떤 논문의 실제 기여가 더 큰 backbone이 아니라 action tokenizer일 수 있고, 다른 논문의 핵심은 data mixture 또는 evaluation protocol일 수 있다.

\section{이 장이 고정하는 것}

이 장은 다음 vocabulary를 고정한다.

\begin{itemize}
  \item Real VLA system: language-conditioned perception을 monitored robot behavior로 바꾸는 전체 책임 구조.
  \item Embodied reasoning layer: task context, constraint, memory, subgoal을 구조화하는 layer.
  \item VLA policy layer: observation, language/subgoal, robot state를 policy-level action $\act_t$로 바꾸는 layer.
  \item Action/control boundary: $\act_t$와 $\ctrlcmd_t$의 구분.
  \item Safety/assurance layer: allow, modify, stop, recover, ask-human 결정을 내리는 layer.
  \item Data/evaluation layer: evidence와 metric, failure taxonomy를 정의하는 layer.
\end{itemize}

후속 장은 이 vocabulary를 바탕으로 전개된다. Chapter 13은 action representation과 control interface를 다룬다. Chapter 14는 robot data engine을 다룬다. Chapter 16은 evaluation protocol을 다룬다. Chapter 19는 safety, assurance, deployment를 다룬다.

\begin{figure}[ht]
\centering
\begin{tikzpicture}[
  node distance=7mm,
  root/.style={draw=rvlaBlue, fill=rvlaBlue!8, rounded corners=1mm, align=center, minimum width=82mm, minimum height=10mm, font=\sffamily\small\bfseries},
  child/.style={draw=rvlaLine, fill=white, rounded corners=1mm, align=left, text width=92mm, minimum height=9mm, font=\sffamily\small},
  arrow/.style={-Latex, thick, draw=rvlaBlue}
]
\node[root] (ch1) {Chapter 1: Real VLA as embodied system stack};
\node[child, below=of ch1] (ch13) {Chapter 13: Action representations and control interfaces\\What is $\act_t$? How does $\act_t$ become $\ctrlcmd_t$?};
\node[child, below=of ch13] (ch14) {Chapter 14: Data engines\\What does a robot trajectory record?};
\node[child, below=of ch14] (ch16) {Chapter 16: Evaluation protocols\\What does benchmark success mean?};
\node[child, below=of ch16] (ch19) {Chapter 19: Safety, assurance, deployment\\Where are the safety gates?};
\draw[arrow] (ch1) -- (ch13);
\draw[arrow] (ch13) -- (ch14);
\draw[arrow] (ch14) -- (ch16);
\draw[arrow] (ch16) -- (ch19);
\end{tikzpicture}
\caption{Chapter 1이 후속 장에 넘기는 dependency map.}
\label{fig:chapter-map}
\end{figure}

\section*{요약}
Real VLA는 하나의 model family가 아니라 책임 구조이다. 핵심은 VLA policy가 내는 $\act_t$를 실제 robot command $\ctrlcmd_t$로 바꾸는 interface, 그 과정에서 필요한 state estimation과 control, safety gate, data/evaluation protocol을 함께 보는 것이다. 이 장의 목적은 최신 system 이름을 나열하는 것이 아니라 후속 장에서 사용할 engineering vocabulary를 고정하는 것이다.

\chapterwork
{OpenVLA, $\pi_0$, FAST, OpenVLA-OFT를 Chapter 1의 responsibility layer 관점에서 다시 읽어라. 모델 이름보다 action/control/data/evidence boundary를 먼저 확인한다.}
{$\act_t$와 $\ctrlcmd_t$를 구분하지 않는 VLA 논문에서는 어떤 failure attribution 문제가 생기는가?}
{mug/rack/glass task에 대해 perception, policy action, controller command, safety monitor, evaluation evidence를 각각 한 문장으로 정의하라.}
